\section{问题定义}
\label{sec:problem}

本节定义 QoS 约束下的逐流选路问题，约定全文使用的数学符号，给出效用函数与优化目标，并在章末论证为何该问题需要用强化学习求解。

\subsection{网络与流量模型}
\label{sec:problem:model}

考虑一个有向网络 $G=(V,E)$，节点集 $V$、链路集 $E$。每条链路 $e\in E$ 具有时延 $d_e$ 与带宽容量 $c_e$；记其在某一时刻的剩余带宽为 $b_e^{\text{res}}\in[0,c_e]$，初始时 $b_e^{\text{res}}=c_e$。

流量以\emph{序列}（episode）为单位在线到达：一个序列内有 $W$ 条流逐条到达，智能体即时为每条流选路、且无法预见后续到达的流。每条流一旦被指派路径，便在本序列内持续占用其带宽——这是在线选路中的”永久连接 / 无限保持时间”假设，本文不建模流的离去。序列结束后容量重置、开始新的一批到达。

第 $t$ 条流由四元组
\begin{equation}
\mathbf{f}_t = (o_t,\, v_t,\, \phi_t,\, b_t),\qquad t=1,\dots,W
\label{eq:flow}
\end{equation}
描述，其中 $o_t,v_t\in V$ 为源宿节点，$b_t$ 为带宽需求，$\phi_t\in[0,1]$ 为 \emph{QoS 意图系数}。$\phi_t$ 服从 Pareto 翻转分布（参数 $\alpha$），重尾集中于低 $\phi$ 区间，使绝大部分流为时延与带宽约束严格的”老鼠流”，仅少量流为约束宽松的”大象流”。带宽需求 $b_t$ 由 $\phi_t$ 通过正相关线性映射加高斯扰动导出：
\begin{equation}
b_t = b_{\text{floor}} + (b_{\text{max}}-b_{\text{floor}})\cdot\phi_t + \mathcal{N}(0,\sigma),
\label{eq:bw_gen}
\end{equation}
使 $\phi_t$ 与 $b_t$ 强正相关（Pearson $r\approx 0.90$）但非完全确定，即 $\phi_t$ 仍携带 $b_t$ 之外的独立 QoS 信息。换言之，$(\phi_t, b_t)$ 联合刻画一条流——大象流”量大且容忍尽力而为”、老鼠流”量小但要求严格保障”，二者构成截然不同的 QoS 画像。流序列从需求分布 $\mathcal{D}$ 中采样，即 $(\mathbf{f}_1,\dots,\mathbf{f}_W)\sim\mathcal{D}$。

处理每条流时，系统为其预先计算 $K$ 条候选最短路径 $\mathcal{P}_t=\{p_t^{1},\dots,p_t^{K}\}$（$K\text{-shortest paths}$，按跳数）。选路决策即从 $\mathcal{P}_t$ 中为 $\mathbf{f}_t$ 指定一条路径 $p_t\in\mathcal{P}_t$ 承载。流被指派后，沿途链路的剩余带宽按实际占用扣减、且在本序列内不再释放，从而改变后续到达流所面对的可用容量——这就是先到流与后到流之间的资源竞争：先到流占用的容量会缩减后到流的可行集。

\subsection{服务质量度量}
\label{sec:problem:metrics}

给定为 $\mathbf{f}_t$ 选定的路径 $p_t$，定义两项服务质量度量。

\textbf{带宽满足比。} 路径可提供的带宽受瓶颈链路约束，记
\begin{equation}
\hat b_t = \min\!\Big(\min_{e\in p_t} b_e^{\text{res}},\; b_t\Big),
\qquad
\rho_t = \hat b_t/b_t,
\label{eq:rho}
\end{equation}
其中 $\rho_t\in[0,1]$（由于 $\hat b_t\le b_t$ 且 $\hat b_t\ge 0$，值域天然有界）为带宽满足比，$\rho_t=1$ 表示需求被完全满足。

\textbf{端到端有效时延。} 我们用一个\emph{拥塞感知的有效时延代理}刻画链路占用对时延的影响：链路占用越满，有效时延越高。记链路占用后的可用率 $u_e = b_e^{\text{res}}/c_e\in[0,1]$，路径有效时延为
\begin{equation}
\hat D_t = \sum_{e\in p_t}\big(1-\log u_e\big)\, d_e,
\label{eq:delay}
\end{equation}
即在物理传播时延 $d_e$ 之上乘以一个随占用上升的拥塞惩罚因子 $(1-\log u_e)$：空载（$u_e\!=\!1$）时因子为 $1$、有效时延回落到物理时延，满载（$u_e\!\to\!0$）时因子发散。需要说明的是，式~\eqref{eq:delay} 并非由排队论严格推导的物理时延，而是一个对数障碍（log-barrier）式的单调拥塞代价——它把策略推离满载链路。在 RL 中它只作为奖励塑形信号，其\emph{单调性}（越拥塞代价越高）而非绝对数值驱动策略学习；相比经典 M/M/1 排队的 $1/(1-\rho)$ 双曲发散，本代理增速更缓、对极端拥塞的惩罚更温和。

\subsection{效用函数与优化目标}
\label{sec:problem:objective}

QoS 意图 $\phi_t$ 决定了一条流”看重什么”。我们用 $\phi_t$ 把带宽与时延两项度量加权成单条流的效用：
\begin{equation}
\begin{split}
U(\mathbf{f}_t, p_t) &=
\underbrace{\big[(1-\phi_t)\,\mathbb{1}[\rho_t\!\ge\!1] + \phi_t\,\rho_t\big]}_{\text{带宽效用 }U^{\text{bw}}_t} \\
&\quad +
\underbrace{(1-\phi_t)\,\frac{1}{1+\hat D_t/\tau}}_{\text{时延效用 }U^{\text{dl}}_t},
\end{split}
\label{eq:utility}
\end{equation}
式~\eqref{eq:utility} 对两类流施加截然不同的目标：

\begin{itemize}
\item 当 $\phi_t\!\to\!0$（老鼠流）：带宽效用退化为硬满足指示 $\mathbb{1}[\rho_t\!\ge\!1]$（仅当带宽被完全满足才得分），且时延效用权重达到最大——即要求严格的带宽保障与低时延。
\item 当 $\phi_t\!\to\!1$（大象流）：带宽效用退化为连续比率 $\rho_t$（部分满足也按比例得分），时延效用项被关闭——即只追求吞吐最大化。
\end{itemize}
其中 $\tau$ 为时延归一化常数（具体取值见第~\ref{sec:exp:hyperparam} 节）。

\textbf{优化目标。} 选路策略需在需求分布 $\mathcal{D}$ 的期望意义下，最大化整个序列的累积效用：
\begin{equation}
\max_{\{p_t\}}\;
\mathbb{E}_{(\mathbf{f}_1,\dots,\mathbf{f}_W)\sim\mathcal{D}}
\left[\sum_{t=1}^{W} U(\mathbf{f}_t, p_t)\right],
\label{eq:objective}
\end{equation}
其中 $\{p_t\}_{t=1}^{W}$ 为各步选路决策构成的策略序列，每步 $p_t\in\mathcal{P}_t$ 依照在线因果约束仅基于当前及历史可观测信息 $\{\mathbf{f}_\tau, p_\tau\}_{\tau\le t}$ 与当前网络状态确定，\emph{不能}依赖尚未到达的未来流 $\{\mathbf{f}_\tau\}_{\tau>t}$。



式~\eqref{eq:objective} 中，给定一条流与一条路径，效用 $U(\mathbf{f}_t,p_t)$ 与剩余带宽的更新都是确定且可精确算出的，这容易让人以为可用经典优化或贪心规则求解。但本问题的真正难点不在记账，而在\emph{未来不可见}：流逐条在线到达，智能体为第 $t$ 条流选路时无法预见 $t+1,\dots,W$ 的后续流（在线因果约束）。若把整条流序列都摆在面前，选路当然可以建成一个整数规划求全局最优；但这样的离线最优解需要预知全部未来，无法在线执行，只能当作事后的性能上界，何况求解本身还是 NP-难的。正因未来不可见，当前每一步选路都会通过剩余带宽 $b_e^{\text{res}}$ 改变后续流的可行集：将老鼠流优先分配至低时延高剩余带宽链路，可能耗尽后来大象流所需的容量；故逐步贪心（如 CSPF，每到一条流即为其选当前满足约束的最短路）必然短视，最优策略必须在对需求分布 $\mathcal{D}$ 的期望意义下、非短视地预留容量并做跨流权衡，即把序列末尾的累积效用回溯归因到早先那几步选路上。这种”在不可见的未来下做序贯决策、并对长程回报做信用分配”的结构，正是强化学习所刻画的马尔可夫决策过程——第~\ref{sec:method} 节据此将式~\eqref{eq:objective} 形式化为 MDP 并给出求解方法。

\label{sec:problem:why_rl}
